\documentclass[../proyecto.tex]{subfiles}
\graphicspath{{\subfix{../images/}}}
\begin{document}

\section{Investigación propuesta}

% máximo diez páginas

\subsection{Presentación del proyecto}

Cyn, no supe qué escribir aquí, para no repetirme.

\subsection{Lugar de enunciación}

Esta práctica artística y académica es una mezcla de confluencias y contradicciones, donde se aspira a no ser aspiracional, a popularizar la producción de ruido para fines artísticos, a promover el goce del sonido desatado en distintas manifestaciones, a desatar modulaciones no hegemónicas, a valorar la infraestructura material y cultural que nos permite sonar y escuchar.

\subsection{Motivaciones}

Las motivaciones que llevan a la realización de esta tesis doctoral son múltiples, entre las que destacan las siguientes.

\subsubsection{Transdisciplina popular artesanal}

La formación para la construcción de sintetizadores generalmente depende de un tránsito entre muchas disciplinas, incluyendo física, ingeniería eléctrica, matemática, carpintería, ciencias de la computación, diseño industrial, bellas artes, entre otros. Esta tesis emerge como una forma de promover que la síntesis de sintetizadores sea una práctica popular transdiciplinar y artesanal \footcite{collins_handmade_2020}, no anecdótica ni alimentada por esfuerzos individuales.

\subsubsection{Materia por sobre cultura}

La investigación es realizada para poder establecer un diálogo con la materia, que emerjan sonidos a partir de componentes que dependen de su fisicalidad, pero no de nociones preconcebidas culturas presentes en la música como armonía, ritmo constante, o afinación.

\subsection{Marco teórico / definiciones conceptuales}

En industria y en español existe muchas veces una amalgama entre conceptos como sintetizadores y teclados, pero en esta tesis nos concentraremos en otros tipos de interfaces \footcite{bjorn_push_2021} y gestualidades para la síntesis sonora.

En la industria contemporánea de la síntesis  también existe mucho énfasis en música, lo que está divorciado del origen artístico e investigativo de muchas personas pioneras del siglo XX que expandieron el campo del arte sonoro y sentaron las bases cartográficas, conceptuales, electrónicas y computacionales \footcite{roads_computer_1996} que permiten la emergencia de un amplio abanico de sintetizadores actuales y futuros.

El mismo concepto de industria está también en tensión con gran parte de las referentes citadas en esta investigación, quienes difuminan las barreras entre luthería, composición e interpretación, muchas veces dedicándose a todas ellas al mismo tiempo.

Cyn, aquí van las definiciones de los conceptos que estoy tratando? Como síntesis?

\subsection{Problema / objeto de estudio / corpus}

El objeto de estudio es la síntesis de sintetizadores, encarnada en empresas que se dedican al diseño, implementación y fabricación de sintetizadores, los libros y materiales educativos que permiten esta práctica transdisciplinar, las instituciones donde se imparten cursos y residencias en torno a la síntesis de sintetizadores, y de forma aún más amplia a la gente desarrolladora de software y hardware y biblitoecas que hace posible la síntesintesis en distintos lenguajes y paradigmas.

\subsection{Criterios de selección del corpus}

En este corpus se incluyen tanto libros de la última década que sistematizan prácticas artísticas y pedagógicas, tanto en electrónica \footcite{pearson_electronic_2024} como en computación \footcite{erwig_erase_2025}.

En el corpus se incluyen también los sintetizadores lanzados por empresas como BASTL Instruments, Ciat-Lonbarde, Critter and Guitari, que han operado bajo principios de software y hardware abierto, que apuestan en la electrónica y la computación como una forma de compartir y de potenciar a artistas a que modifiquen sus propios instrumentos.

\subsection{Referencias críticas sobre objeto de estudio / estado del arte}

A nivel de crítica sobre los objetos de estudio, se referenciarán los estudios de Yuk Hui \footcite{hui_arte_2025}, y los manifiestos y propuestas de permacomputación, donde se proponen alternativas anticapitalista de consumo de electrónica y computación, reciclando materiales existentes, permitiendo la reparabilidad, y la legalidad de las licencias que artistas y sintesintesistas puedan compartir sin pasar a llevar sus creaciones \footcite{tozzi_history_2017}.

\subsection{Referentes artísticos y contexto en que se sitúa la práctica artística}

Los grandes referentes artísticos sobre los que se sitúa esta práctica artística incluyen a los siguientes gigantes de la sintesíntesis 

\subsubsection{Ciat-Lonbarde}

Peter Blasser es un artista que ha desarrollado muchas líneas de sintetizadores en los últimos 20 años, entre ellos los samplers, órganos y máquinas de ritmo de Ciat-Lonbarde, los instrumentos programables de Shbobo como Shnth y Shtar, y su más reciente familia de sintetizadores Cafeteria Quantum.

\subsubsection{Clacktronics}

La empresa Clacktronics de Inglaterra lanzó un libro con paneles prepicados para romperlos y luego poblarlos con electrónica y convertirlos en un sintetizador Eurorack, incluyendo la carcasa.

\subsubsection{Corazón de robota}

La artista Constanza Piña ha desarrollado una práctica de creación de sintetizadores, material de enseñanza DIY, talleres e infraestructura para enseñar y promover la creación de sintetizadores a partir de circuitos integrados y distintos sensores.

\subsubsection{Critter and Guitari}

Actualmente esta empresa tiene 2 productos principales: un sintetizador de sonido programable en el lenguaje Pure Data, y un sintetizador de video programable en el lenguaje Python.

\subsubsection{KORG}

La empresa japonesa KORG ha lanzado en la última década una línea de sintetizadores programables con su propia SDK, incluyendo a minilogue xd, drumlogue, prologue y nts-1 mk1 y mk2.

\subsubsection{Oficina de sonido}

La empresa de Daniel Llermaly ha lanzado varios sintetizadores, entre ellos un sistema modular con el estándar Eurorack.

\subsubsection{Sonic Liberation Devices}

La práctica de esta cooperativa está situada en torno a la computación aplicada a la música como una forma de resistencia política, además de estar enlazada tanto el diseño y la manufactura de sintetizadores con su enseñanza y el compartir.

\subsection{Novedad de la propuesta}

El estado del arte ha mostrado que existen muchos referentes en el mundo que pueden ser catalogados como popusintesíntesis en muchos de sus preceptos.

Gran parte del corpus presupone distintos tipos de conocimientos para públicos masivos, desde estar en inglés, tener entrenamiento en electrónica o en desarrollo de software, mecánica, arte conceptual, o incluso música.

Esta propuesta entonces es novedosa al hacerse cargo de la dimensión metafórica y conceptual necesaria para entender las artes sonoras, de tener un enfoque primero conceptual y no enfocado en la implementación o en algún software o hardware en particular, y de tener como precepto la realización de herramientas paramétricas que permitan no tener un sintetizador en particular, sino que una constelación de subproductos y artefactos que permiten la emergencia de nuevas síntesis. También se incluye la dimensión de reflexión de infraestructura tanto comercial como académica que permite que esta investigación florezca.

\end{document}